\section{配合比模型与目标函数}
\label{sec:model}
\subsection{决策变量与材料用量}
以 $1\,\mathrm{m^3}$ 模型拌合物为计算基准，取
\begin{equation}
  \boldsymbol{x}=(B,s,r,P)^{\mathsf T},\qquad
  \begin{cases}
    300\leq B\leq460,\\
    0\leq s\leq0.55,\\
    0.32\leq r\leq0.52,\\
    0\leq P\leq15.
  \end{cases}
  \label{eq:bounds}
\end{equation}
$B$ 为胶凝材料总质量，$s$ 为其中矿渣的质量比例，$r$ 为有效水胶比，
$P$ 为处理后干意面的质量。$B$ 和 $P$ 的单位均为 $\mathrm{kg/m^3}$。
这些区间用于构造有界演示问题，不是规范限值或工程建议。
意面不计入胶凝材料，因而增加 $P$ 不会改变水胶比分母。
水泥、矿渣与总加水量分别为
\begin{equation}
  C=(1-s)B,\qquad S=sB,\qquad W=rB+aP,\qquad a=0.8.
  \label{eq:mass}
\end{equation}
这里 $aP$ 是简化的补偿吸水量，$a$ 的单位为 $\mathrm{kg/kg}$；
假设这部分水被添加物吸收，不参与有效水胶比计算。该设定只用于区分总水与有效水，
未描述吸水动力学、释水过程或体积膨胀。

固定含气体积为 $0.02\,\mathrm{m^3}$，粗骨料体积为 $0.38\,\mathrm{m^3}$，
则粗骨料质量 $G=1026\,\mathrm{kg}$。余下体积由细砂补足：
\begin{equation}
  A=2650\left(1-0.02-0.38-
  \frac{C}{3150}-\frac{S}{2900}-\frac{W}{1000}-\frac{P}{1350}\right).
  \label{eq:volume}
\end{equation}
各分母为演示密度，单位为 $\mathrm{kg/m^3}$。上述绝对体积关系同时计入干添加物和补偿水，
采用组分体积可加的理想化近似。其意义在于保证比较对象具有相同模型体积，
避免通过少配材料使碳排与成本同时下降。真实材料的表观密度、孔隙水和干湿状态
需要采用相互一致的试验定义，不能直接沿用本式中的数值。

\subsection{碳排放与成本核算}
令材料索引 $i\in\{C,S,W,A,G,P\}$，其质量、单位碳排因子和单位价格依次为
$m_i,e_i,c_i$。定义模型边界内的碳排放和成本：
\begin{equation}
  E(\boldsymbol{x})=\sum_i e_im_i,\qquad
  C_{\mathrm{cost}}(\boldsymbol{x})=\sum_i c_im_i.
  \label{eq:environment-cost}
\end{equation}
$E$ 的单位为 $\mathrm{kg\,CO_2e/m^3}$，$C_{\mathrm{cost}}$ 的单位为元/$\mathrm{m^3}$。
\tabref{tab:factors}中的所有数值均由演示情景设定，未从文献或供应商数据库中标定。

\begin{table}[htbp]
  \centering
  \caption{材料核算参数（全部为合成情景设定）}
  \label{tab:factors}
  \begin{tabular}{@{}lrrr@{}}
    \toprule
    材料 & 密度/(kg/m$^3$) & $e_i$/(kg CO$_2$e/kg) & $c_i$/(元/kg)\\
    \midrule
    水泥 & 3150 & 0.86000 & 0.500\\
    矿渣 & 2900 & 0.12000 & 0.300\\
    水 & 1000 & 0.00035 & 0.004\\
    细砂 & 2650 & 0.00500 & 0.090\\
    粗骨料 & 2700 & 0.00400 & 0.075\\
    处理后干意面 & 1350 & 0.10000 & 0.060\\
    \bottomrule
  \end{tabular}
\end{table}

对意面采用简化的废弃物切断分配：食品原生产负担不进入本情景，收集、清洗、干燥、粉碎等
处理负担由 $e_P$ 统一代表，不另给负的处置抵扣。由此得到的结论依赖这一分配设定。
若改用经济分配、质量分配，或把运输与避免处置收益纳入边界，必须重算目标函数。
混凝土环境评价的功能单位与系统边界会影响比较结论\citep{heede2012}，
因此这里的 $E$ 仅表示限定边界内的演示指标，不能解释为完整生命周期排放。

该模型也不预设添加意面必然减碳。在其他变量固定时，增加一单位 $P$ 会减少细砂质量
$2650(0.8/1000+1/1350)$，但需要增加补偿水和处理负担。基准情景下有
\begin{equation}
  \frac{\partial E}{\partial P}
  =0.1+0.00035\times0.8
   -0.005\times2650\left(\frac{0.8}{1000}+\frac{1}{1350}\right)>0.
  \label{eq:pasta-carbon}
\end{equation}
因而废弃物再利用与降低本边界碳排并不等价；二者是否相容，应由完整的目标与约束判断。

\subsection{假设性能响应与可行域}
本案例未提供强度、流变或耐久性试验。为使优化问题可重复运行，直接定义下列合成响应，
不把它们称为经验拟合、试验预测模型或经过验证的材料定律：
\begin{align}
  F(\boldsymbol{x})&=68+0.045(B-360)-120(r-0.38)
                    -20s-P-0.03P^2,\label{eq:strength}\\
  H(\boldsymbol{x})&=160+800(r-0.38)-60s-1.5P
                    +0.1(B-360).\label{eq:slump}
\end{align}
$F$ 的单位为 MPa，$H$ 的单位为 mm，各系数的量纲由相应变量确定。
模型将较高有效水胶比设定为强度代价，将添加物设定为强度与流动性的负担，
意在构造需要取舍的优化情景，而不是证明这些作用的真实方向或幅度。
特别是矿渣在不同龄期、养护条件下的作用并不能由式\eqref{eq:strength}概括。

除变量边界外，规定
\begin{equation}
  F\geq40,\quad 140\leq H\leq220,\quad
  140\leq W\leq220,\quad 500\leq A\leq1000.
  \label{eq:constraints}
\end{equation}
强度下限 $40\,\mathrm{MPa}$ 是本演示另行选取的阈值，与题目中的案例编号“42”无关。
砂量约束排除不合理的体积分配，水量与坍落度约束限制计算配合比的搜索区域。
这些条件只定义数学可行性，满足它们不代表材料通过结构安全、耐久性或环境适用性验证。

为了消除不同约束违反量的量纲差异，使用按预设尺度归一化的总违反度
\begin{equation}
  \begin{aligned}
  v(\boldsymbol{x})={}&\frac{[40-F]_+}{40}
  +\frac{[140-H]_++[H-220]_+}{80}\\
  &+\frac{[140-W]_++[W-220]_+}{80}
  +\frac{[500-A]_++[A-1000]_+}{500},
  \end{aligned}
  \label{eq:violation}
\end{equation}
其中 $[z]_+=\max(z,0)$。变量边界由遗传算子的有界实现保证，
式\eqref{eq:violation}用于对其余约束进行排序。程序保留原始性能值及违反度，
不把不可行候选伪装成低碳可行解。

\subsection{多目标优化问题}
优化问题写为
\begin{equation}
  \min_{\boldsymbol{x}\in\Omega}
  \boldsymbol{f}(\boldsymbol{x})=
  \bigl(E(\boldsymbol{x}),C_{\mathrm{cost}}(\boldsymbol{x}),-F(\boldsymbol{x})\bigr),
  \label{eq:multiobjective}
\end{equation}
其中 $\Omega$ 由式\eqref{eq:bounds}与式\eqref{eq:constraints}共同确定。
将强度写成负值，是为统一最小化方向。成本与碳排可能同向变化，
但它们并非相同指标；即使二者高度相关，也不应在没有分析的情况下删除其中一个。
强度目标使模型不仅寻找刚好达到阈值的方案，也保留以更高资源投入换取强度余量的选择。
由于性能函数并无实测支撑，本文不将最大化 $F$ 解释为提高真实服役可靠度。

\section{NSGA-II求解与方案选择}
\label{sec:algorithm}
\subsection{约束支配与非支配排序}
若两个候选均可行，当 $\boldsymbol{f}(\boldsymbol{x})$ 在所有分量上不劣于
$\boldsymbol{f}(\boldsymbol{y})$，且至少一个分量严格更小时，称 $\boldsymbol{x}$ 支配
$\boldsymbol{y}$。含约束时采用三条规则：可行解优于不可行解；
两个不可行解中总违反度较小者优先；两个可行解按普通 Pareto 支配关系比较。
违反度相同的不可行候选之间不再用目标函数强行构造支配关系。

NSGA-II通过非支配排序得到逐层前沿，再通过拥挤距离保留目标空间中的分散性\citep{deb2002}。
在同一前沿中，对第 $j$ 个目标排序，内部候选的拥挤距离增量为
\begin{equation}
  d_i^{(j)}=
  \frac{f_j^{(i+1)}-f_j^{(i-1)}}{f_j^{\max}-f_j^{\min}}.
  \label{eq:crowding}
\end{equation}
边界候选优先保留；若某目标在该前沿内为常数，则跳过该目标的归一化项，避免除零。
拥挤距离只在同一等级内用于保留多样性，并不表示候选的工程优劣。

\subsection{遗传算子与精英保留}
每代从当前种群进行二元锦标赛：先比较非支配等级，等级相同再比较拥挤距离，
仍相同则随机打破平局。连续变量采用有界模拟二进制交叉和多项式变异。
交叉产生的新值限制在变量区间内，变异对每个变量独立作用。
父代与子代合并后形成两倍规模的候选集，按前沿逐层装入下一代；
最后一层装不下时，按拥挤距离由大到小截取。
该步骤实现精英保留，但不意味着有限代数下找到了全局 Pareto 前沿。

\begin{algorithm}[htbp]
  \caption{带约束支配的 NSGA-II 配合比优化}
  \label{alg:nsga}
  \begin{algorithmic}[1]
    \Require 变量区间、合成响应、核算因子、种群数 $N$、迭代数 $T$、随机种子
    \Ensure 最终可行非支配集合及间隔记录
    \State 在变量区间内采样 $N$ 个候选，评价并记录初始状态
    \For{$t=1,\ldots,T$}
      \State 计算父代的非支配等级与拥挤距离
      \State 进行锦标赛选择、有界交叉和变异，产生 $N$ 个子代
      \State 评价子代并与父代合并
      \State 按约束支配关系分层，按等级和拥挤距离保留 $N$ 个候选
      \State 每隔 10 代及末代记录可行候选数和最低碳排
    \EndFor
    \State 提取最终种群内的可行非支配候选，去除重复解
    \State \Return 候选集、运行参数和计算记录
  \end{algorithmic}
\end{algorithm}

本文采用统一评价预算比较优化算法和随机搜索。每次优化包括初始种群和 $T$ 代子代，
总评价次数为 $N(T+1)$；随机基线在相同变量边界内评价同样数量的候选，
随后保留其可行非支配集合。评价次数统计包含不可行候选，避免因只统计成功样本而偏向某个方法。
多次运行使用独立种子，报告差异而不只挑选最有利的一次。

\subsection{折中方案与比较口径}
在主运行得到的可行非支配集合 $\mathcal{P}$ 上，对三目标分别作极差归一化：
\begin{equation}
  z_j(\boldsymbol{x})=
  \frac{f_j(\boldsymbol{x})-\min_{\boldsymbol{y}\in\mathcal{P}}f_j(\boldsymbol{y})}
       {\max_{\boldsymbol{y}\in\mathcal{P}}f_j(\boldsymbol{y})-
        \min_{\boldsymbol{y}\in\mathcal{P}}f_j(\boldsymbol{y})}.
  \label{eq:normalize}
\end{equation}
若分母为零，该分量置零。取等权理想点距离最小的候选为展示用折中方案：
\begin{equation}
  \boldsymbol{x}^{\star}=\arg\min_{\boldsymbol{x}\in\mathcal{P}}
  \left[\frac{1}{3}\sum_{j=1}^{3}z_j(\boldsymbol{x})^2\right]^{1/2}.
  \label{eq:compromise}
\end{equation}
这一规则只是显式的偏好设定，不是唯一最优决策。极差取决于当前候选集合，
新候选进入后，折中方案可能发生改变。因此同时保留最低碳排、最低成本和最高合成强度方案，
并公开完整前沿，使读者可以用其他权重重算选择。

参考方案固定为 $(B,s,r,P)=(380,0.15,0.42,0)$，用于说明变化方向。
相对碳排变化按 $(E_{\mathrm{ref}}-E)/E_{\mathrm{ref}}$ 计算，
成本同理；强度另外列出绝对数值和变化量，不把强度损失隐藏在减碳百分比中。
该参考配合比本身也是演示设定，并不代表现行工业平均水平。
